Since PDFs are non-perturbative objects, encoding low-scale physics, they are not computable from first principles inside pQCD.
Instead, PDFs are obtained from a fitting procedure, which can be summarized as follows:
\begin{conclusionbox}{PDF fitting}
    \begin{enumerate}
        \item Guess a candidate PDF
        \item Compute the theory predictions with this candidate PDF
        \item Compare the theory predictions to the experimental measurement
    \end{enumerate}
    Repeat the above steps until the procedure converges, i.e.\ the comparison is \enquote{good enough}.
\end{conclusionbox}
The specific details on how to guess a candidate PDF, how to exactly compute the theory predictions, which experimental data to use, and what \enquote{good enough} exactly means is a choice.
That is the reason why competing PDF fitting groups exist and why they do not necessarily obtain the same PDF, although we all assume the latter to be universal.

The main point for us here is that we need to compute theory predictions very many times for slightly different PDFs.
This is a serious problem for complicated hadronic observables, such as, e.g., jet cross sections~\cite{NNLOJET:2025rno}, where the computation can take many hours.
The solution is to compute fast interpolation grids, using, e.g., PineAPPL~\cite{Carrazza:2020gss,christopher_schwan_2025_15635174}, APPLgrid~\cite{Carli:2010rw}, or fastNLO~\cite{Wobisch:2011ij}, which store the partonic matrix elements in an efficient way, such that any subsequent convolution with an arbitrary PDF is fast.

In DIS this is not a very big problem, since typically coefficient functions, \cref{eq:fact1}, are known analytically and thus the convolution with PDFs is straightforward.
However, PineAPPL allows also DIS cross sections to be stored thus we can use this simplified case to introduce the general concept here.
Moreover, having all observables in a PDF fit sharing the same technology also simplifies their consistent treatment and inside the NNPDF collaboration this is realized in the pineline framework~\cite{Barontini:2023vmr}.

Interpolation grids are only a discretized version of the coefficient functions which we can recombine in a suitable way a posteriori.
Specifically, they are higher dimensional tensors on which we perform fast linear algebra operations.
The first three major dimensions are:
\begin{description}
    \item[bin] This is the list of cross section that we need to compute simultaneously.
        In DIS this typically corresponds to all $N_\text{dat}$ data point measured in a given configuration, e.g.\ in \cref{tab:HERA}.
        Thus, in practice the actual list is given by the experimental measurement.
    \item[order] This is the list of perturbative orders, i.e.\ we store each $\vb C^{(k)}(z)$, \cref{eq:CpQCD}, separately.
        Having this list available allows to use a given grid also at lower orders and thus study, e.g., perturbative stability.
        In practice, the actual list is given by the requested accuracy of the user (provided we can actual fulfill that request).
        Recall that LO, \cref{eq:LOC}, and NLO, e.g.\ \cref{eq:NLOCq}, are quite different and this generates correlations with the momentum fraction dimension (see below).
    \item[channel] This is the list of PDF flavor combinations, i.e.\ the sum over $j$ in \cref{eq:fact1}.
        Having this list available allows to study the impact of different PDF combination and, e.g., determine which PDF the observable can potentially constrain most.
        In DIS it is given by the coupling structure of the exchanged boson, i.e.\ $\gamma/Z$ vs.\ $W^\pm$.
        In practice, the actual list is the easiest given directly in terms of individual parton flavors as the available combinations are too complex (see \cref{sec:sv} for additional complications).
\end{description}

These three dimensions are present in any interpolation grid and in DIS there are two additional dimensions:
\begin{description}
    \item[scale] This is the list of factorization and renormalization scales, discussed in \cref{sec:sv}.
        In DIS this dimension turns out to be trivial, i.e.\ it has support by a single value, as in fully inclusive DIS there is only a single scale, $Q^2$, available and we set typically $\mu_F^2 = \mu_R^2 = Q^2$.
        In practice, the actual list is given by the $Q^2$ specified in the bin configuration.
    \item[momentum fraction] This is the list of momentum fractions $z$, \cref{eq:fact1}, over which the coefficient functions are convolved with the PDF.
        In DIS we are facing the problem of distributions, see \cref{sec:rsl}, i.e.\ we can not directly tabulate the coefficient functions themselves.
        In practice, the actual list is given by a convolution with a suitable set of interpolation functions, which form a basis for PDFs.
\end{description}

For fully inclusive DIS we end up with five dimensional tensors, where each inner dimension may depend on the outer.

\fherror{explain interpolation?}
